\chapter{Lezione 17/03}

Il motivo per cui introduciamo questo concetto di sostituzione è perché il lemma o teorema che poi introdurremo legato a questo concetto di sostituzione ci permetterà di ottenere formule valide a partire da altre formule valide.

Questo sarà una prima forma di deduzione logica molto elementare, ma che ci introdurrà verso calcoli deduttivi più raffinati e complessi.

\begin{theorem}[namedlabel={thm:substitution}{Teorema di Sostituzione}]{Teorema di Sostituzione}
  Sia \(\phi\in TAUT\) e \(\sigma\) un'arbitraria sostituzione. Allora \(\phi\sigma\in TAUT\).
\end{theorem}

\begin{note}{}
  Una formulazione alternativa e più intuitiva di questo teorema è che la sostituzione preserva la tautologicità (o la validità) delle formule.
\end{note}

Per dimostrare questo teorema, andremo a introdurre una versione semantica della sostituzione.

Dati \(\phi, \sigma\) e \(\alpha\) (assegnamento per \(\phi\sigma\)), definiamo \(\alpha^\sigma:AP(\phi)\to\set{0,1}\) come segue:
\[\alpha^\sigma(P)=val^\alpha(\sigma(P))\]

Possiamo vedere l'applicazione di \(\sigma\) come una mera sostituzione sintattica (renaming), mentre questa nuova definizione ci dà una sostituziuone semantica.

Vediamo dunque un risultato intermedio che ci sarà utile per dimostrare il \Namedref{thm:substitution}.
\begin{lemma}[namedlabel={lem:substitution}{Lemma di Sostituzione}]{Lemma di Sostituzione}
  \(\forall \phi \in PL, \sigma \in PL^{AP}, \alpha \in ASS(\phi\sigma)\) vale che
  \[val^\alpha(\phi\sigma)=val^{\alpha^\sigma}(\phi)\]
\end{lemma}


\begin{proof}
  Dimostriamo per induzione strutturale.

  Nel caso base \(\phi \in AP\). Allora \(\phi\sigma = \sigma(\phi)\).

  Ma quindi
  \[val^{\alpha}(\phi\sigma) = val^\alpha(\sigma(\phi)) = \alpha^\sigma(\phi) \overset{def}{=} val^{\alpha^\sigma}(\phi)\]

  Andando quindi ai casi induttivi, iniziamo con \(\phi = \neg \psi\). In questo caso dalla definizione di sostituzione abbiamo
  \[\phi\sigma = (\neg\psi)\sigma = \neg(\psi\sigma)\]

  Allora:
  \[val^{\alpha}(\phi\sigma)= val^{\alpha}(\neg(\psi\sigma)) = 1 - val^{\alpha}(\psi\sigma)\]
  Ma per ipotesi induttiva \(val^{\alpha}(\psi\sigma) = val^{\alpha^\sigma}(\psi)\), quindi
  \[val^{\alpha}(\phi\sigma)= 1 - val^{\alpha^\sigma}(\psi) = val^{\alpha^\sigma}(\neg\psi) = val^{\alpha^\sigma}(\phi)\]

  Per gli altri casi induttivi la struttura di dimostrazione è analoga, e si basa sempre sull'ipotesi induttiva.
\end{proof}

Andiamo quindi alla dimostrazione del \Namedref{thm:substitution}.
\begin{proof}{\Namedref{thm:substitution}}
  Vogliamo dimostrare che \(\forall \sigma \in PL^{AP}\), \(\phi \in TAUT \implies \phi\sigma \in TAUT\).

  Procediamo per contrapposizione, ovvero assumiamo che \(\phi\sigma \not\in TAUT\) e dimostriamo che \(\phi \not\in TAUT\).

  Se \(\phi\sigma \not\in TAUT\) allora \(\exists \alpha \in Ass(\phi\sigma)\st val^\alpha(\phi\sigma) = 0\).

  Ma per il \Namedref{lem:substitution}, \(val^\alpha(\phi\sigma) = val^{\alpha^\sigma}(\phi)\), quindi \(val^{\alpha^\sigma}(\phi) = 0\), e quindi \(\phi \not\in TAUT\).
\end{proof}

\begin{observation}{}
  È importante notare che nonostante la sostituzione preserva validità e tautologicità, essa non preserva la soddisfacibilità o la verità in generale.

  Formalmente:
  \[\phi \in Sat \not\implies \phi\sigma \in Sat\]

  Questo è facile da vedere considerando \(\phi = P \in AP\). Si ha che \(\phi \in Sat\) banalmente da qualsiasi assegnamento \(\alpha\) tale che \(\alpha(P)=1\).
  Scegliendo però \(\sigma : P \to P\land \neg P\), che è chiaramente insoddisfacibile.
  Di conseguenza, fissato un modello \(m\), si ha che
  \[m\models\phi \not\implies m\models\phi\sigma\]
  Altrimenti otterei una contraddizione.

  Cosa diversa per la insoddisfacibilità. Infatti
  \[\phi \in Unsat \iff \neg \phi \in TAUT \implies (\neg\phi)\in TAUT\iff \]
  \[\iff \neg(\phi\sigma) \in TAUT \implies \phi\sigma \in Unsat\]
\end{observation}

Possiamo notare che \(A \to (B\to C)\equiv (A\land B)\to C\). Chiamiamo la prima formula \(\phi_1\) e la seconda \(\phi_2\).

Sappiamo che \(\phi_1\) può essere falsa se e soltanto se \(A\) è vera, \(B\) è vera e \(C\) è falsa. 
Invece essendo \(\phi_2\) una semplice implicazione questa è falsa se e soltanto se \(A\land B\) è vera e \(C\) è falsa, ovvero se e soltanto se \(A\) è vera, \(B\) è vera e \(C\) è falsa. Quindi \(\phi_1\) è falsa se e soltanto se \(\phi_2\) è falsa, e quindi \(\phi_1\equiv\phi_2\).

Possiamo trasformare questa equivalenza in una tautologia, scrivendo \(A \to (B\to C) \leftrightarrow (A\land B)\to C\).

\begin{theorem}[namedlabel={thm:semantic_deduction}{Teorema di Deduzione Semantica}]{Teorema di Deduzione Semantica}
  Sia \(\Gamma \subseteq PL\) e \(\phi,\psi \in PL\). Allora
  \[\Gamma,\phi \models \psi \iff \Gamma \models \phi\to\psi\]
\end{theorem}

\begin{proof}
  Dalla definizione di conseguenza logica, \(\Gamma,\phi \models \psi\) significa che
  \[\forall m \in 2^{AP}, m\models \Gamma \cup \set{\phi} \implies m\models \phi\]

  L'argomento sinistro dell'implicazione è però equivalente a dire che \(m\models \Gamma\) e \(m\models \phi\). Quindi
  \[\forall m \in 2^{AP}, m\models \Gamma \text{e} m\models \phi \implies m\models \psi\]

  Dando quindi semantica agli operatori di metalivello nel modo usuale (ovvero identico a quelli della logica proposizionale), e possiamo riscrivere la precedente formula con variabili proposizionali \(A\), \(B\) e \(C\) al posto di \(m\models \Gamma\), \(m\models \phi\) e \(m\models \psi\) rispettivamente, ottenendo
  \[\forall m \in 2^{AP}, A\text{ e }B \implies C\]

  Ma dalla tautologia vista in precedenza, abbiamo che
  \[\forall m \in 2^{AP}, m\models \Gamma \implies (m\models \phi \implies m\models \psi)\]

  Per semantica dell'implicazione però \(m\models \phi \implies m\models \psi \iff m\models \phi \to \psi\), dunque
  \[\forall m \in 2^{AP}, m\models \Gamma \implies m\models \phi\to\psi\]
\end{proof}

\section{Forme normali}

Abbiamo un linguaggio che contiene molti simboli, ma abbiamo già osservato tramite alcune tatutologie che è possibile esprimere alcuni operatori in funzione di altri.

Un esempio potrebbero essere le leggi di De Morgan
\[\neg(\phi\land\psi) \equiv \neg\phi \lor \neg\psi\]
\[\neg(\phi\lor\psi) \equiv \neg\phi \land \neg\psi\]

Interpretando queste equivalenze come regole di trasformazione, questo ci permette di spingere le negazioni fino alle proposizioni atomiche mantenedo l'equivalenza logica.

In questo modo, possiamo definire una prima forma normale, chiamata PNF (Positive Normal Form) o NNF (Negation Normal Form), che contiene tutte le formule di PL in cui le negazioni occorrono solo davanti a proposizioni atomiche.

Per definire formalmente questa forma normale, introduciamo il concetto di letterale

\begin{definition}{Letterale}
  Una formula \(\phi\) è un letterale se \(\phi \in AP\) o \(\phi = \neg P \), con \(P \in AP\).
\end{definition}

\begin{definition}{Forma normale positiva}
  Definiamo la forma normale positiva (PNF o NNF) come l'insieme di tutte le formule costruite come in PL, senza usare negazioni e utilizzando l'insieme dei letterali come AP
\end{definition}

Si ha ovviamente che \(PNF \subset PL\), ma dalle leggi di De Morgan è ovvio che \(\forall \phi \in PL \exists \phi' \in PNF: \phi \iff \phi'\)

In generale il concetto di forma normale è un sott'insieme di formule costruibile in maniera induttiva tale che mantenga completezza semantica, ovvero che per ogni formula esista una formula equivalente in forma normale.

Le due forme normali più importanti sono la forma normale congiuntiva (CNF) e la forma normale disgiuntiva (DNF).

\begin{definition}{Disjuntiva Normal Form}
  Definiamo la forma normale disgiuntiva (DNF) come l'insieme di tutte le formule con la seguente struttura:
  \[\phi = \bigvee_{i=1}^k(\bigwedge_{j=1}^{r_i} l_{ij})\]
  dove \(l_{ij}\) è un letterale per ogni \(i\) e \(j\).
\end{definition}

\begin{example}{}
  Un esempio di formula in DNF è la seguente:
  \[(P\land \neg Q) \lor (\neg P \land \neg Q \land R) \lor \ldots\]

  Ognuna delle parentesi (quindi ogni congiunzione) è detta clausola congiuntiva.
\end{example}

\begin{definition}{Conjunctive Normal Form}
  Definiamo la forma normale congiuntiva (CNF) come l'insieme di tutte le formule con la seguente struttura:
  \[\phi = \bigwedge_{i=1}^k(\bigvee_{j=1}^{r_i} l_{ij})\]
  dove \(l_{ij}\) è un letterale per ogni \(i\) e \(j\).
\end{definition}

\begin{example}{}
  Un esempio di formula in CNF è la seguente:
  \[(P\lor \neg Q) \land (\neg P \lor \neg Q \lor R) \land \ldots\]

  Ognuna delle parentesi (quindi ogni disgiunzione) è detta clausola disgiuntiva.
\end{example}

Essendoci dualità tra congiunzione e disgiunzione, è facile vedere che le due forme normali sono duali.

Ovviamente \(CNF \cup DNF \subseteq PNF\).

\todo{Esempio tabella verità.}

\chapter{Lezione 19/03}

Per la \(\CNF\) il ragionamento è simile, poiché dobbiamo esprimere la formula come una congiunzione di clausole disgiuntive, che cozza con l'uso degli assegnamenti come in precedenza, poiché gli assegnamenti sono delle congiunzioni.

Quindi invece di costruire una formula elencando tutti i casi in cui è vera ma costruiremo una formula che non è vera in tutti gli assegnamenti in cui quella originale era falsa.

Quindi basta creare una formula in cui ogni clausola è della forma forma analoga a quella della \(\DNF\) ma con una negazione davanti, tutte queste congiunte. Ma dalla leggi di de morgan queste clausole congiuntive sono trasformabili in clausole disgiuntive.

\todo{fai esempio pratico}

Se sono tutti a \(1\) posso prendere una qualsiasi tautologia che è già in \(\CNF\) come \(P \lor \neg P\) e mappare in questa.

Queste costruizioni non sono di certo ottimali, poiché per ogni \(\phi\) la formula \(\phi'\) sarà molto più grande. Infatti la grandezza della tabella di verità è \(2^{\AP{\phi}}\), e quindi nel peggiore dei casi si può avere un numero esponenziale sulle proposizioni atomiche di clausole. Non è detto che sia sempre così ma ci sono alcune formule in cui è inevitabile, in particolare quando si passa dall'una all'altra.

\begin{example}{}
  Sia \(\phi \in\DNF\) tale che \(\phi\) sia:
  \[(P_1 \land P_2)\lor (P_3 \land P_4)\]

  Andiamo quindi a trasformarla in \(\CNF\) usando le tautologie note
  \[(P_1 \land P_2)\lor (P_3 \land P_4)\equiv\]
  \[\equiv ((P_1 \land P_2)\lor P_3)\land((P_1 \lor P_4)\lor P_4)\equiv\]
  \[\equiv (P_1 \lor P_3) \land (P_2 \lor P_3) \land (P_1 \lor P_4) \land (P_2 \lor P_4)\]

  Si è passati quindi da \(2\) a \(4\), che può sembrare un incremento lineare. Possiamo notare però che questo processo può essere vista come una funzione di scelta.

  In particolare, data
  \[\phi = (l_1^1 \land l_2^1) \lor (l_1^2 \land l_2^2) \lor\ldots\lor (l_1^n \land l_2^n)\]
  Possiamo costruire una classe di funzione \(f: \set{1,\ldots,n}\to \set{1,2}\) tale che ogni clausola della trasformazione corrisponde a una di queste funzioni di scelta, dunque avremo \(\abs{\phi'}\leq 2^{k}\).

  SPIEGA MEGLIO
\end{example}

Il motivo per cui abbiamo introdotto queste forme è perché se la formula è in \(\DNF\) sarà molto computazionalmente facile valutare al sua soddisfacibilità, mentre se è in \(\CNF\) sarà molto facile valutare la sua validità.

Infatti sia \(\phi = C_1 \lor C_2 \lor\ldots\lor C_k\), con \(C_i = (l_1 \land \ldots \land l_n), \forall i \in \set{1,\ldots, k}\). Una formula del genere è soddisfacibile se e solo se almeno una delle \(C_i\) è soddisfacibile, ma ogni \(C_i\) è insoddisfacibile se e solo se esistono due letterali \(l_j\) e \(l_{j'}\) tali che \(l_j = \neg l_{j'}\). Questo perché se ci sono tali letterali, la clausola contiene la contraddizione \(P \land \neg P\), altrimenti è sembre possibile costruire un assegnamento che rende vera la clausola, in maniera analoga ma inversa a come abbiamo ragionato precedentemene con le tabelle di verità.

Dualmentne, se \(\phi = D_1 \land D_2 \land\ldots\land D_k\), con \(D_i = (l_1 \lor \ldots \lor l_n), \forall i \in \set{1,\ldots, k}\). Una formula del genere è valida se e solo se tutte le \(D_i\) sono valide, ma ogni \(D_i\) è valida se e solo se esistono due letterali \(l_j\) e \(l_{j'}\) tali che \(l_j = \neg l_{j'}\).

\section{Calcolo deduttivo NUOVO CAPITOLO}

Abbiamo definito una logica come una tripla \((\mathcal{L}, \models, \vdash)\), dove \(\mathcal{L}\) è un linguaggio, \(\models\) è la relazione di soddisfacibilità che ne esprime la semantica, e \(\vdash\) è un calcolo, che ci permette di dedurre formule a partire da altre, o verificare se una formula è valida o meno.

Abbiamo già visto una formula di calcolo tramite le tabelle di verità, ma questo tipo di calcolo non \qi{scala}, poiché ad esempio nel prim'ordine non sarà possibile ridursi a un numero finito di casi, come abbiamo visto per il prim'ordine.

Il calcolo che vedremo sarà di pura manipolazione simbolica, noi ne vedremo due, il primo è il \keyword{calcolo di deduzione naturale}.

Un calcolo in generale è definito da un insieme di formule denominato \keyword{assiomi} \(Axi \subseteq PL\),  e un insieme di regole di trasformazione dette \keyword{regole di inferenza} \(R\) che a partire da un numero finito di formule, dette \keyword{premesse}, ci permettono di dedurre una nuova formula, detta \keyword{conclusione}.

Ovviamente non tutte le regole di trasformazioni sono regole d'inferenza valide, poiché vorremmo che queste regole preservassero delle proprietà semantiche. Tramite esse potremmo definire la relazione di \keyword{deducibilità} \(\Gamma \vdash \phi\), se partendo da \(\Gamma\) applicando regole in \(R\) si potrà dedurre \(\phi\). L'obiettivo sarebbe fare in modo in maniera puramente sintattico per ottenere la conseguenza logica, ovvero vorremmo che:
\[\Gamma \models \phi \iff \Gamma \vdash \phi\]

Per quanto riguarda gli assiomi, sono un insieme minimale di tautologie dalle quali tramite il teorema di sostituzione varranno in diverse istansiazioni.\footnote{In un certo senso la regola di sostituzione fa parte di \(R\).}

Il primo calcolo che vedremo però non ha assiomi, ma solo regole di inferenza, e si chiama \keyword{calcolo di deduzione naturale}.

Ovviamente per definirlo è sufficiente definire le regole, il cui numero dipende dal numero di simboli logici che abbiamo. In particolare avremo regole associate ad ogni operatore logico, e verrà usato \(\bot \equiv P\land \neg P\) come simbolo di contraddizione. Le regole catturano il significato semantico degli operatori logici.

\todo{Inserisci regole di deduzione naturale}

AND
\(A \land B \vdash A\) \(A \land B \vdash B\) \(A , B \vdash A\land B\)
IMPLIES
\(A \to B, A \vdash B\) \(B \vdash A \to B\), questa regola mi permette di eliminare un assunzione, poiché se da \(A\) deduco \(B\), allora posso dedurre \(A \to B\) senza più assumere \(A\).
OR
\(A \vdash A \lor B\) \(B \vdash A \lor B\) \([A]\ldots C ,[B]\ldots C, A \lor B \vdash C\) 
NOT
\(\neg A , A \vdash \bot\) \([A]\ldots \bot \vdash \neg A\)
BOT
\(\bot \vdash A\) \([A] \ldots \bot \vdash A\) (RAA)

Per il processo di deduzione l'idea è che partiremo da delle assunzioni e usando passi atomici di deduzione, che sono applicazioni di regole di inferenza, arriveremo a dedurre la formula che ci interessa. Alcune delle assunzioni potranno essere eliminate, tutte le altre saranno invece assunzioni residue, che andranno a formare \(\Gamma\) nella deducibilità \(\Gamma \vdash \phi\).